\chapter{جداول الصفحات والذاكرة الافتراضية}
\label{chap:mem}

تُعتبر جداول الصفحات الآلية الأكثر مرونة وقوة التي يقدمها العتاد لنظام التشغيل لإدارة الذاكرة.
تتيح \indextext{جداول الصفحات (Page Tables)} للنواة ربط العناوين الافتراضية للعمليات بالعناوين الفيزيائية للحاسوب،
مما يضمن عزلاً محكماً لمساحات الذاكرة بين العمليات المختلفة، ويوفر حماية كاملة لذاكرة النواة.

\section{عتاد إدارة الذاكرة بـ RISC-V}

تعتمد معمارية RISC-V بنية العناوين الشجرية ثلاثية المستويات والمعروفة بـ \indextext{Sv39}؛
حيث تبلغ مساحة العنوان الافتراضي 39 بت.
تتكون العناوين من رقم الصفحة الافتراضية (Page Number) المكون من 27 بت، والإنزياح (Offset) المكون من 12 بت بداخل الصفحة (بحجم صفحة فيزيائية يبلغ 4096 بايت).

تحفظ النواة العنوان الفيزيائي لجذر جدول الصفحات في السجل الخاص \texttt{satp}.
وعند كل عملية وصول للذاكرة، يقرأ عتاد وحدة إدارة الذاكرة (MMU) السجل \texttt{satp} ويتصفح المستويات الثلاثة لجدول الصفحات للوصول إلى \indextext{مدخلة جدول الصفحة (Page Table Entry - PTE)}.

تضم مدخلة جدول الصفحة PTE العنوان الفيزيائي ومجموعة من بتات الصلاحيات والحماية:
\begin{itemize}
\item B بت الصلاحية \texttt{PTE\_V}: يحدد ما إذا كانت المدخلة صالحة وموجودة.
\item B بت القراءة \texttt{PTE\_R}: يتيح قراءة البيانات من الصفحة.
\item B بت الكتابة \texttt{PTE\_W}: يتيح كتابة البيانات بداخل الصفحة.
\item B بت التنفيذ \texttt{PTE\_X}: يتيح للمعالج تنفيذ التعليمات الموجودة بالصفحة.
\item B بت المستخدم \texttt{PTE\_U}: يحدد ما إذا كان يُسمح لبرامج وضع المستخدم الوصول للصفحة.
\end{itemize}

لتسريع عملية ترجمة العناوين وتجنب القراءة المتكررة للذاكرة، يتضمن المعالج تسريعاً عتادياً يُسمى «مخزن الترجمة المؤقت (Translation Lookaside Buffer، TLB)».
يحفظ مخزن الترجمة المؤقت نتائج الترجمة الأخيرة للعناوين؛ فتتحقق إصابة في مخزن الترجمة المؤقت (TLB Hit) لتوفير الوقت، أو يحدث إخفاق في مخزن الترجمة المؤقت (TLB Miss) فيتصفح المعالج جدول الصفحات.

\section{توزيع الذاكرة في \texttt{xv6}}

تُقسم النواة مساحة العنوان الفيزيائي لـ RISC-V إلى منطقتين أساسيتين عند العنوان \texttt{0x80000000}:
المنطقة الأولى دون هذا العنوان تُخصص لأجهزة الإدخال والإخراج المربوطة بالذاكرة (MMIO)،
بينما تُخصص المنطقة ابتداءً من \texttt{0x80000000} لذاكرة RAM الفيزيائية التي تضم شفرة النواة ومساحات التخصيص الديناميكي.

تخصص النواة لكل عملية مستخدم جدول صفحات مستقل تماماً يُحمل في السجل \texttt{satp} عند تنقيل السياق إليها.
تبدأ مساحة عنوان المستخدم من العنوان الصفر \texttt{0x0} وتنمو صعوداً عبر استدعاءات التخصيص \texttt{sbrk}.

\section{إدارة الذاكرة الفيزيائية}

يتولى الموزع الفيزيائي المسمى \texttt{kalloc.c} إدارة صفحات الذاكرة الحرّة المتاحة بداخل النواة.
ينظم الموزع الصفحات الفيزيائية المتاحة بداخل \indextext{قائمة خطية ترابطية (Free List)}؛
فعند طلب صفحة عبر \texttt{kalloc()}، تُقطع الصفحة الأولى من القائمة وتُصفر بياناتها؛
وعند تحرير صفحة عبر \texttt{kfree(pa)}، تُعاد الصفحة إلى أسر القائمة الترابطية.
